\section{Extensions and Discussion}

\begin{frame}{Extension 1: Intermediate Goods}
  \small
  Many sectors sell output to firms as well as households. The extension lets
  goods be used as intermediate inputs.

  \vspace{0.10cm}
  \begin{block}{Modified production function}
    \begin{equation*}
      F^i=A_i n_i k_i^\alpha q_i^\beta,\qquad
      \alpha,\beta>0,\quad \alpha+\beta<1 .
    \end{equation*}
  \end{block}

  \begin{alertblock}{Balanced-growth restriction}
    The intermediate-input aggregator must be Cobb-Douglas:
    \[
      \eta=1.
    \]
  \end{alertblock}

  \vspace{0.10cm}
  \slidecite{The paper notes that U.S. input-output data show large intermediate-use shares across sectors.}
\end{frame}

\begin{frame}{Extension 2: Many Capital Goods}
  \small
  Instead of a single manufacturing capital good, suppose there are $\kappa$
  capital-producing sectors aggregated into $G$.

  \vspace{0.12cm}
  \begin{block}{Relative employment across capital-goods sectors}
    \begin{equation*}
      \frac{n_{m_j}}{n_{m_i}}
      =
      \left(\frac{\xi_{m_j}}{\xi_{m_i}}\right)^\mu
      \left(\frac{A_{m_i}}{A_{m_j}}\right)^{1-\mu}.
    \end{equation*}
    \begin{equation*}
      \frac{\dot n_{m_j}}{n_{m_j}}-\frac{\dot n_{m_i}}{n_{m_i}}
      =(1-\mu)(\gamma_{m_i}-\gamma_{m_j}).
    \end{equation*}
  \end{block}

  \vspace{0.08cm}
  The logic mirrors the baseline model, but the relevant substitution elasticity
  is now $\mu$ inside the capital-goods aggregator.
\end{frame}

\begin{frame}{Many Capital Goods: Balanced-Growth Implication}
  \small
  \begin{alertblock}{Key result}
    Multiple capital-producing sectors with different TFP growth rates can
    coexist on an aggregate balanced growth path only if
    \[
      \mu=1.
    \]
  \end{alertblock}

  \vspace{0.10cm}
  \begin{itemize}
    \item The aggregate capital-sector growth rate is a weighted average of
      capital-sector TFP growth rates.
    \item If $\mu\neq1$, the weights depend on relative productivity levels and
      therefore change over time.
    \item If $\mu=1$, the Cobb-Douglas aggregator fixes expenditure/input shares,
      so the aggregate growth rate is constant.
  \end{itemize}
\end{frame}

\begin{frame}{Strengths of the Paper}
  \small
  \begin{columns}[T,onlytextwidth]
    \begin{column}{0.48\textwidth}
      \begin{block}{Theoretical strengths}
        \begin{itemize}
          \item Clean separation between preferences and technologies.
          \item Structural change generated without nonhomothetic preferences.
          \item Aggregate dynamics reduce to a Ramsey benchmark under $\theta=1$.
        \end{itemize}
      \end{block}
    \end{column}
    \begin{column}{0.48\textwidth}
      \begin{alertblock}{Interpretive strengths}
        \begin{itemize}
          \item Gives a precise version of Baumol's mechanism.
          \item Explains why relative prices and employment can rise together.
          \item Reconciles large nominal shifts with small real-share shifts.
        \end{itemize}
      \end{alertblock}
    \end{column}
  \end{columns}
\end{frame}

\begin{frame}{Limitations and Discussion Questions}
  \small
  \begin{itemize}
    \item \focus{No nonhomotheticity}: the model intentionally abstracts from
      income-elasticity channels that are important in other structural-change models.
    \item \focus{Identical production shapes}: sectoral differences enter mainly
      through TFP growth, not through capital intensity or fixed factors.
    \item \focus{Balanced-growth knife edge}: exact coexistence requires $\theta=1$,
      though the paper argues deviations may be quantitatively small.
    \item \focus{Sector classification}: real industries often produce final,
      intermediate, and capital goods at the same time.
  \end{itemize}

  \vspace{0.20cm}
  \begin{alertblock}{Discussion prompt}
    Is the low final-goods substitution elasticity more plausible at broad
    sectoral levels than at detailed industry levels?
  \end{alertblock}
\end{frame}

\begin{frame}{Where the Literature Goes Next}
  \scriptsize
  \begin{alertblock}{Big question after Ngai and Pissarides (2007)}
    Is structural change driven mainly by relative-price effects, or do we also
    need income effects, capital deepening, trade, and input-output linkages?
  \end{alertblock}

  \vspace{0.05cm}
  \begin{tabularx}{\textwidth}{@{}p{0.24\textwidth}p{0.30\textwidth}X@{}}
    \toprule
    Research direction & Representative papers & Main idea \\
    \midrule
    Price effects vs. income effects
      & Herrendorf, Rogerson, and Valentinyi (2013, 2014)
      & Compare CES price-effect models with nonhomothetic-preference models
        in U.S. structural-transformation data. \\
    Flexible nonhomothetic preferences
      & Boppart (2014); Alder, Boppart, and M\"uller (2022);
        Comin, Lashkari, and Mestieri (2021)
      & Reintroduce income effects to explain Engel curves, falling agriculture,
        rising services, and hump-shaped manufacturing. \\
    Other supply-side mechanisms
      & Acemoglu and Guerrieri (2008); Alvarez-Cuadrado et al. (2017, 2018)
      & Emphasize capital deepening, sectoral capital intensity, and differences
        in capital-labor substitution. \\
    Quantitative and open-economy extensions
      & Swiecki (2017); Uy, Yi, and Zhang (2013); Sposi (2019)
      & Quantify competing mechanisms and study how trade and input-output
        linkages shape structural change. \\
    \bottomrule
  \end{tabularx}
\end{frame}

\begin{frame}{Main Takeaways}
  \small
  \begin{enumerate}
    \item Unequal sectoral TFP growth changes relative prices.
    \item CES demand converts relative price changes into expenditure-share changes.
    \item If $\varepsilon<1$, labor moves toward sectors with lower TFP growth.
    \item Aggregate balanced growth can coexist with this movement if $\theta=1$.
    \item Intermediate inputs and multiple capital goods preserve the core idea
      under Cobb-Douglas aggregation restrictions.
  \end{enumerate}

  \vspace{0.25cm}
  \begin{alertblock}{Bottom line}
    Structural change can be a permanent feature of sectoral allocation while
    the aggregate economy still looks balanced.
  \end{alertblock}
\end{frame}
